% =============================================================================
% Supplementary material for the SphereQuant paper.
% Drops in as an appendix after the main paper, or as a separate file with
% \appendix preceding it.
% =============================================================================

\section{Detailed analysis of the architectural boundary}
\label{app:boundary}

Section~\ref{sec:experiments-boundary} of the main text reports that
SphereQuant trails QuaRot-RTN on MobileNet-V2 at four-bit and six-bit, and on
EfficientNet-B0 at six-bit and eight-bit. This appendix decomposes the
loss into three distinct mechanisms, each predicted by a separate aspect
of the methodology in Section~\ref{sec:method}, and identifies a
practical recipe for closing the gap.

\subsection{Mechanism 1: small-fan-in failure of the Beta approximation}
\label{app:beta-failure}

The codebook used by SphereQuant is the Lloyd--Max minimum-MSE quantizer
for $\mathrm{Beta}(d/2, d/2)$ on $[-1, 1]$. The standard deviation of
this distribution scales as $(d+1)^{-1/2}$, giving the values listed in
Table~\ref{tab:fanin-beta} of the main text. The reasoning that
motivated the codebook (Section~\ref{sec:method-motivation}) was that a
sharply concentrated post-rotation density admits levels that cluster
near the origin and outperforms a uniform grid in proportion to that
concentration. When the density itself is wide, the same reasoning
predicts that the matched codebook collapses onto a near-uniform
alternative. Concretely, at $d = 9$ the Lloyd--Max levels at four bits
deviate from evenly spaced positions by less than five percent, so the
codebook offers essentially no advantage over QuaRot's uniform grid.

A separate sample-complexity issue compounds this. At $d = 2048$ the
empirical histogram of one row's coordinates contains 2048 samples and
already approximates the asymptotic $\mathrm{Beta}(1024, 1024)$
density. At $d = 9$ each row contains nine samples; even if the
underlying distribution were exactly Beta, no codebook can be matched
to a nine-sample empirical density. Pooling coordinates across rows
would defeat the per-row L2 normalization that anchors the unit-sphere
geometry.

This mechanism is the geometric content of
Figure~\ref{fig:rotation-vision} in the main text. The leftmost column
of the MobileNet-V2 and EfficientNet-B0 rows shows the post-rotation
histogram remaining nearly flat across $[-1, 1]$, with the predicted
Beta density itself nearly flat and offering no concentration to
exploit.

\subsection{Mechanism 2: L2-versus-absmax normalization on outlier rows}
\label{app:l2-vs-absmax}

For a single row $w \in \mathbb{R}^d$, SphereQuant divides by
$\|w\|_2 = \sqrt{\sum_i w_i^2}$ before quantizing. RTN-absmax and
QuaRot-RTN divide by $\|w\|_\infty = \max_i |w_i|$. The two norms are
related by $\|w\|_\infty \in [\|w\|_2 / \sqrt{d},\, \|w\|_2]$, with the
upper bound attained when one coordinate holds all of $w$'s energy.

The relative behavior of the two normalizations on a $3 \times 3$
depthwise kernel illustrates the effect. Suppose eight of the nine
weights have magnitude $\epsilon$ and the ninth has magnitude $1$. Then
$\|w\|_\infty = 1$ and $\|w\|_2 = \sqrt{1 + 8\epsilon^2} \approx 1$.
Under L2 normalization, the outlier maps near $1$, in the sparsest
region of the Beta codebook (where centroid spacing approaches the
maximum at any bit width), and the per-coordinate error from rounding
the outlier dominates the row's quantization error. Under absmax
normalization the outlier maps exactly to the integer level
$L = 2^{b-1} - 1$ and is reconstructed without rounding error at any
$b \geq 2$, while the eight near-zero coordinates compress to small
magnitudes that the uniform grid handles coarsely but symmetrically.

The L2-versus-absmax tradeoff is therefore a separate axis from the
Beta-concentration argument: even if Mechanism~\ref{app:beta-failure}
were absent, an outlier-dominated row would still favor absmax
quantization over L2 normalization. Depthwise convolutions concentrate
their energy in one or two coordinates because the kernel is small;
large-fan-in linear layers spread their energy across many coordinates
and therefore favor L2 normalization. The two mechanisms reinforce
each other on the same set of layers, which is why the boundary is so
sharp.

\subsection{Mechanism 3: regime change as the bit budget grows}
\label{app:regime-change}

A mixed architecture such as EfficientNet-B0 contains layers in three
fan-in regimes simultaneously: pointwise convolutions in the
$d \geq 96$ regime, depthwise convolutions at $d \in \{9, 25\}$, and
squeeze-excitation bottlenecks at $d \in \{4, 8, 16, 32\}$. The
total quantization-induced accuracy drop is approximately the sum of
per-layer contributions, weighted by each layer's downstream
sensitivity.

As the bit budget grows from four to six to eight, all three regimes
move toward zero error, but they do so at different rates. The
large-$d$ regime saturates first, because both rotation methods reach
near-FP32 reconstruction on those layers as soon as the per-coordinate
error falls below the per-row signal-to-noise threshold. The small-$d$
regimes saturate later, because the codebook mismatch in
Mechanism~\ref{app:beta-failure} is bit-rate-independent: a flat Beta
density does not become more concentrated when more bits are
available. QuaRot-RTN, lacking this bit-rate-independent penalty,
catches up to FP32 on the small-$d$ layers faster.

The empirical signature of this regime change is visible in the
EfficientNet-B0 numbers reported in
Section~\ref{sec:experiments-vision}. At four-bit, where the
pointwise-convolution error dominates, SphereQuant outperforms
QuaRot-RTN by 31.2 percentage points. At six-bit, where the pointwise
convolutions are saturated and the depthwise plus squeeze-excitation
error dominates, QuaRot-RTN wins by 2.5 percentage points. At
eight-bit the gap narrows to 0.7 points in favor of QuaRot-RTN as
both methods approach FP32. The same regime change explains the
absence of an analogous reversal on TinyLlama and Phi-1.5: every
linear layer in both models sits in the favorable regime, so there is
no second regime to come dominant as the budget grows.

\subsection{A practical layer-adaptive dispatch}
\label{app:hybrid}

The three mechanisms above suggest a simple deployment recipe.
Before quantizing, inspect each layer's fan-in $d$. Apply SphereQuant
when $d \geq 32$ and fall back to QuaRot-RTN when $d < 32$. The
threshold is chosen at the boundary where the Beta standard deviation
on $[-1, 1]$ exceeds $0.17$ and the empirical advantage over a uniform
grid disappears. The dispatch rule itself requires no calibration
data, only the layer's fan-in. We expect this hybrid to recover the
architectural-boundary losses on MobileNet-V2 and EfficientNet-B0
without sacrificing the gains in the large-$d$ regime that SphereQuant
provides on every other architecture in our study, but a complete
evaluation across all six bit widths and all six architectures is
left to future work.

\subsection{Empirical Lloyd--Max as a fallback codebook}
\label{app:empirical-lloyd}

A second avenue for closing the boundary gap is to replace the
asymptotic Beta codebook with an empirical Lloyd--Max codebook fit to
the actual post-rotation coordinate histogram of each layer. The
empirical codebook adapts to small-$d$ departures from the Beta limit,
recovers the Lloyd--Max optimality guarantee for whatever the layer's
post-rotation density actually is, and remains training-free in the
sense that the codebook is constructed entirely from the model's own
weights. The codebook construction is one-time per layer: collect the
flattened post-rotation coordinates of all rows, run a few iterations
of Lloyd--Max with $2^b$ levels, and store the resulting centroids
alongside the codes. Storage cost grows by one centroid table per
unique layer rather than one per unique fan-in, which is bounded above
by the number of layers in the model and is negligible compared to
the weight tensor.

We do not report a full evaluation of the empirical Lloyd--Max
codebook in this paper because the asymptotic Beta codebook is
sufficient for the architectures where SphereQuant is most useful. The
empirical alternative is the natural construction for a follow-up
study targeting the depthwise-separable regime.

\subsection{Summary}

The architectural boundary that limits SphereQuant on MobileNet-V2 and
on EfficientNet-B0 at higher bit widths is fully characterized by the
three mechanisms above and is consistent with the geometric argument
in Section~\ref{sec:method-boundary}. The boundary is not a defect of
the method; it is the prediction of the same mathematics that
motivates the construction. Two adaptations,
Sections~\ref{app:hybrid} and \ref{app:empirical-lloyd}, would close
the gap without sacrificing the wide-fan-in gains, at the cost of
small additional bookkeeping. We leave their empirical evaluation to
future work and include the present analysis to make the boundary
behavior fully explicable from the geometry alone.
